\chapter{Gestión de errores en C-Horizon}
\label{cap:errores}

Para cada error se indicará su tipo y su posición en el archivo (fila y columna). La clase GestionErroresExp se
encarga de actuar conforme a cada error. Además, llevará un contador de errores que imprimirá cuando haya habido
algún error. Cuando no, permitirá la generación de código. Vamos a explicar los distintos tipos de errores. 

\section{Errores léxicos}
Simplemente se imprimirá el token no reconocido.\\

Input [input6\_errlexico.txt]\footnote[1]{Cuando pongamos un input para mostrar la salida, indicaremos así el nombre del archivo
de texto concreto para poder ejecutar el código. Estos archivos se encuentran en la carpeta tests.}
\begin{algorithm}	
	\Integer{}(\Integer{}) main(a) \{ \\
	\quad \Integer{} hola ::= 2; \comm{//Error léxico} \\
	\quad \Whilee{} (\const{True}) \{ \\
	\quad\quad hola\operator{++};\_\comm{//Error léxico} \\
	\quad\} \\
	\quad	\Returnn a\;
	\}\;
\end{algorithm}

Output

\begin{itemize}
	\item \err{[Fila 2, Columna 11] Caracter inesperado: ":"}
	\item \err{[Fila 4, Columna 10] Caracter inesperado: "\_"}
\end{itemize}
\newpage
\section{Errores sintácticos}
Si el error se encuentra en una instrucción simple, se retomará la 
compilación justo donde termina la instrucción, esto es, después de ";". En caso de que el error se dé en la declaración de un ámbito (una función, una 
condición, un bucle) se retomará la compilación justo al terminar ese ámbito. Es decir, se buscará el caracter $"\{"$ y se omitirá todo el código hasta 
encontrar el "\}" que lo cierre (para ello también se llevará un contador de los "\{" sin cerrar). Algunos ejemplos de errores sintácticos: 

Input [input7\_errsintaxis.txt]
\begin{algorithm}
	\Integer{}(\Integer{}) main(a) \{ \\
	\quad \Integer{} hola := 2\;
	\quad \Whilee{} (\Integer{}) \{ \comm{//Error sintáctico en la condición}\\
	\quad \quad hola\operator{++}\;
	\quad \} \\
	\quad	\Returnn \operator{.};\comm{//Error sintáctico en la condición}\\
	\}\;
\end{algorithm}

Output

\begin{itemize}
	\item \err{[Fila 3, Columna 9] Elemento inesperado ``int''}
	\item \err{[Fila 6, Columna 9] Elemento inesperado ``.''}
\end{itemize}
\newpage
\section{Errores de vinculación}

Todos los nodos del árbol AST implementan dos funciones: bindFunc y bind. La primera, bindFunc,
solo inserta una entrada en la tabla de símbolos actual con la declaración de una función. Esto lo
hacemos para permitir que una función declarada en un bloque pueda ser usada en cualquier punto del
bloque. Por lo tanto, cada vez que abramos un nuevo bloque, añadiremos una nueva tabla de símbolos
para después llamar siempre primero a bindFunc y luego a bind del prgrama de ese bloque. Al terminar
de vincular el bloque desapilaremos la tabla de símbolos. Además al desapilarla, guardaremos la tabla
del bloque porque nos es útil luego (para la definición de structs y clases y para el tipado). Las tablas 
están implementadas como mapas de clave String y valor Declaracion. En vez de tener una pila de tablas,
hemos usado una pila de pilas de tablas. Hay varias razones. AQUÍ EXPLÍCALO 

Cuando lleguemos, usando bind, a una declaración de variable, insertaremos dicha declaración en el tabla de 
símbolos actual. Al insertar una declaración de variable comprobaremos que no se ha insertado otra declaración
de variable con el mismo identificador. Sí se puede declarar una variable y una función con el mismo identificador.
Como también existe sobrecarga, en vinculación permitimos que dos declaraciones de funciones tengan mismo
identificador. Será en tipado donde contrastaremos si ambas tienen mismo tipo funcional (y lanzaremos un error
de ambigüedad).

Al usar un identificador,se buscará en hacia atrás en la pila de tablas y se añadirán al conjunto de vínculos del
identificador todas las declaraciones que encajen. Cuando en un bloque se encuentre una declaración de variable que
coincida, no se buscará más declaraciones de variables en los bloques previos (ocultación). Sin embargo, sí se seguirán buscando
declaraciones de funciones a las que poder vincular el identificador.

Por último, también llevaremos paralelamente otra pila de tablas que recoja los tipos nuevos (clases o structs) declarados en cada
bloque. Cuando se defina un tipo nuevo, insertaremos en la tabla de tipos actual una entrada nueva. Esta tabla funciona
exactamente igual que la anterior pero los pares son String (nombre del tipo nuevo) y NodoTipo (el tipo concreto). 
\newpage
Vamos a ver un ejemplo de un código que usa lo que acabamos de explicar y no da ningún error.

Input [input13\_errvinc.txt]

\begin{algorithm}
	\Integer{}(\Integer{}) f (a) \{\}\;
	\Boolean{} f\;
	\Integer{}() main()\{\\
		\quad \Integer{} f; \comm{//Declaración de variable "f" que no produce error por ocultación}\\
		\quad \Integer{}(\Boolean{}) f (a) \{\}\;
		\quad \comm{//Lo siguiente no produce error de declaración ni de tipado porque f se ha vinculado correctamente a la declaración del primer bloque}\\
		\quad \Integer{} x := f(2)\;
		\quad \comm{//Lo siguiente tampoco proudcirá error de tipado porque f ha sido vinculada correctamente a la declaración de este bloque}\\
		\quad x := f\;
	\};
\end{algorithm}

A continuación, veremos los tipos de errores concretos de vinculación. \\ \\ \\ \\
\newpage
\subsection{Identificador ya declarado}

Este error se lanza cuando se declara un identificador que ya ha sido declarado. \\

Input [input8\_errvinc.txt]

\begin{algorithm}
	\Integer{} hola\;
	\Boolean{} hola; \comm{//Error por identificador ya declarado}\\
	\Integer{}(\Integer{}) hola (a) {}\;
	\Integer{}(\Integer{}) hola (a) {}; \comm{//Resulta en un error de tipado, no de vinculación}\\
	\Void () main () \{\};
\end{algorithm}

Output\footnote{Solo hemos incluido los errores de vinculación aquí}

\begin{itemize}
	\item \err{[Fila 2, Columna 6] Identificador ya declarado: hola}
\end{itemize}

\subsection{Identificador no declarado}

Este error se lanza cuando utiliza un identificador que no ha sido declarado prevaimente. \\

Input [input9\_errvinc.txt]

\begin{algorithm}
	\Integer{}() main() \{\\
	\quad \Integer{} x\;
	\quad x := a \operator{+} 1; \comm{//Error de vinculación porque no se ha declarado ``a''} \\
	\quad x := f(a); \comm{//Solo hay error de id no declarado en a} \\

	\quad \{\\
	\quad \quad \Integer{} a\;
	\quad \quad x := a \operator{+} 1; \comm{//Esta instrucción sí que es correcta} \\
	\quad \}

	\quad x := a \operator{+} 1; \comm{//Ahora vuelve a dar error porque se ha cerrado el bloque} \\
	\};\\
	\Integer{}(\Integer{}) f (num) {}\;
\end{algorithm}

Output

\begin{itemize}
	\item \err{[Fila 3, Columna 6] Identificador no declarado: a}
	\item \err{[Fila 4, Columna 8] Identificador no declarado: a}
	\item \err{[Fila 9, Columna 6] Identificador no declarado: a}
\end{itemize}

\subsection{Tipo nuevo ya definido}

Este error se lanza cuando se define un tipo nuevo con un nombre con el que ya se había definido
un tipo nuevo antes. Aunque sean tipos nuevos distintos (una clase y un struct), si ambos tienen el mismo
nombre lanzamos error porque puede producir ambigüedad en el futuro. \\

Input [input10\_errvinc.txt]

\begin{algorithm}
	\Class A \{\\
		\quad \Integer{} num\;
	\};\\
	\Struct A \{\};  \comm{//Error porque ya se ha definido un tipo nuevo con el nombre ``A''}\\
	\Class A a;  \comm{//Esta declaración es correcta, pues solo se acepta la primera definición de ``A''}\\
	\Integer{}() main() \{\\
	\quad a.num := 2\;
	\};
\end{algorithm}

Output

\begin{itemize}
	\item \err{[Fila 4, Columna 8] Tipo nuevo ya definido con ese nombre: A}
\end{itemize}

\subsection{Tipo nuevo no definido}

Este error se lanza cuando se declara una variable de un tipo nuevo que no ha sido declarado previamente. \\

Input [input11\_errvinc.txt]

\begin{algorithm}
	\Struct A a; \comm{//Error porque no se ha definido ese tipo nuevo}\\
	\Struct A \{\}; \comm{//Ahora sí se ha declarado, las declaraciones no darán error}\\
	\Struct A a1\; 
	\Integer{} () main () \{\\
		\quad \Struct A a2\;
		\quad \Struct B \{\}\;
	\}\;

	\Struct B b; \comm{//Error porque no hay una declaración del struct B visible}\\
\end{algorithm}

Output

\begin{itemize}
	\item \err{[Fila 1, Columna 1] Tipo nuevo no definido: A}
	\item \err{[Fila 8, Columna 1] Tipo nuevo no definido: B}
\end{itemize}

\subsection{Error en el número de argumentos}

Este error se lanza cuando se declara una función y el número de argmuentos no corresponde con el número de los
tipos de argumentos de la función. Este error lo detectamos en vinculación porque cuando se vincula la cabecera de 
la función, se insertan declaraciones en la tabla de símbolos con cada uno de los argumentos, para que los identificadores
del ámbito de la función se puedan vincular a estas declaraciones. Para que estas declaraciones no choquen con otras
declaraciones de ese bloque, al vincular los argumentos de una función abrimos un nuevo bloque intermedio entre el bloque
donde está definida la función y el bloque del ámbito de la función.\\

Input [input12\_errvinc.txt]

\begin{algorithm}
	\Integer{}(\Integer{}) main (a,b) \{\};
\end{algorithm}

Output

\begin{itemize}
	\item \err{[Fila 1, Columna 10] Error en el número de argumentos de la función: f}
\end{itemize}

\subsection{Error por falta de main}
Si no se declara en el ámbito global una función main, se lanzará un error. Debe ser en el ámbito global.
La sintaxis impide que se pueda declarar en cualquier otro ámbito.\\

Input [input16\_errNoMain.txt]

\begin{algorithm}
	\comm{//No hay función main}\\
	\Integer{}(\Integer{}, \Integer{}) suma (a,b) \{\\
		\quad \Returnn a \operator{+} b\;
	\}\;
	\Integer{}(\Integer{}) factorial (n) \{\\
		\quad \Ifff{(n \operator{==} 0)} \{ \\
		\quad \quad \Returnn 1\;
		\quad \}\\
		\quad \Returnn n \operator{*} factorial(n\operator{-}1)\;
	\}\;
\end{algorithm}

Output

\begin{itemize}
	\item \err{[ERROR] No se ha declarado la función "main"}
\end{itemize}

\newpage
\subsection{Error por varios main}
Se han detectado varios main, lo que impide la generación de código.

\begin{algorithm}
	\comm{//Error por haber varios main}\\
	\Void () main() \{\}\;
	\Integer{}(\Integer{}) main(arg)\{\}\;
\end{algorithm}

Output

\begin{itemize}
	\item \err{[ERROR] Se han declarado varias funciones "main"}
\end{itemize}


\newpage
\section{Errores de tipado}

El tipado lo hemos gestionado de la siguiente manera. En primer lugar, hemos creado más tipos que los tipos usuales
de nuestro lenguaje de programación (enteros, booleanos, struct, etc). Estos tipos sirven para dos cosas: detectar errores
y tipar instrucciones. El tipado de instrucciones sirve para localizar instrucciones imposibles en ciertos contextos. Por
ejemplo, una declaración de función en un struct o una instrucción \Return{} en una función void, entre otros.

Cada nodo recibe un conjunto de tipos esperados. Y, conforme a ese conjunto, se tipa él mismo y a sus hijos. Si el tipo del
nodo no coincide con ninguno de los tipos esperados, se lanza error de tipado mostrando los tipos que se esperaban. Repetimos
que estos tipos también pueden ser tipos de instrucciones, como DECVARIABLE (una declaración de variable) entre otras. Después 
de tiparse, cada nodo devuelve hacia arriba un conjunto con los posibles tipos válidos para ese nodo.

Vamos a ver algunos errores a continuación. No enseñaremos todos pues son muchos. Al final de esta sección se mencionarán
unos cuantos ficheros de texto preparados con los errores señalados con comentarios.
\\

Input [input14\_errtipado.txt]

\begin{algorithm}
	\Integer{}() main() \{\\
		\quad \Integer{}(\Integer{}) f(a)\{\}\;
		\quad \Boolean{}(\Integer{}) f(a)\{\}\;
		\quad f(2); \comm{//Error por ser una expresión ambigua}\\
		
		\quad \Void () f () \{\\
		\quad 	\quad return \const{True}; \comm{//Error, las funciones void no devuelven nada}\\
		\quad \}\;
		\quad \Integer{} x\;
		\quad \comm{//En los structs solo se pueden declarar variables (inicializadas o no) y definir otros tipos nuevos}\\
		\quad \Struct A  \{ \\
		\quad 	\quad \Integer{} a := 2; \comm{//Correcto}\\
		\quad 	\quad \Struct B \{\}\;
		\quad 	\quad \Void() f2() {}; \comm{//Error, no se puede declarar una función en un struct}\\
		\quad 	\quad x := 2; \comm{//Error, no se puede realizar una asignación (si no es una declaración) en un struct}\\
		\quad \}\;	
	\};
\end{algorithm}
\newpage

Output\footnote{Solo hemos incluido los errores de tipado}

\begin{itemize}
	\item \err{[Fila 4, Columna 2] Expresión ambigua}
	\item \err{[Fila 6, Columna 3] Error de tipado. Se esperaba uno de los siguientes tipos: DEFSTRUCT DEFCLASE        INSTRUCCIÓN     ASIGNACIÓN      DECVARIABLE     DECFUNCION}
	\item \err{[Fila 13, Columna 3] Instrucción no válida}
	\item \err{[Fila 14, Columna 3] Instrucción no válida}
\end{itemize}

Es en el tipado donde se resuelve definitivamente la vinculación de identificadores. Se quitan del conjunto de vínculos aquellas declaraciones
que no coincidan con ningún tipo esperado. Si queda más de un vínculo, lanzaremos un error de expresión ambigua. Si no queda ningún vínculo,
lanzaremos también error de tipado. Este es un ejemplo sencillo.
\\

Input [input15\_errtipado.txt]

\begin{algorithm}
	\comm{//Estas dos funciones son distintas porque su tipo de devolución es distinto}\\
	\Boolean{}(\Integer{}) f(a)\{\}\;
	\Integer{}(\Integer{}) f(a)\{\}\;
	
	\Integer{}() main() \{ \\
		\quad f(2);\comm{ //Error de ambigüedad, pues se espera cualquier tipo de instrucción}\\
		\quad \Integer{} x := f(2); \comm{//Aquí se resuelve y no hay error de ambigüedad, pues solo encaja la segunda función (porque el tipo de retorno es el que se espera en la asignación)}\\
		\quad \Integer{}[2] arr := f(2); \comm{//Error de tipado porque ninguna declaración del identificador ``f'' encaja con los tipos esperados}\\
	\};
\end{algorithm}

Output



\begin{itemize}
	\item \err{[Fila 5, Columna 2] Expresión ambiguaa}
	\item \err{[Fila 7, Columna 16] Error de tipado. Se esperaba uno de los siguientes tipos: ENTERO[2](ENTERO)}
\end{itemize}

El primer error muestra que, al ser tomada la expresión funcional como una instrucción, no importa el tipo de retorno, y 
el compilador no puede distinguir entonces entre las dos declaraciones de "f", así que se produce un error de ambigüedad.

El segundo error muestra que se esperaba un tipo funcional con tipo de retorno un array de entero de dos dimensiones
y un parámetro de tipo entero. Esto es, porque el tipo de retorno era el tipo esperado que venía determinado por la 
asignación y el tipo de los parámetros por los parámetros de esa expresión funcional. Sin embargo, no existe ninguna 
declaración de ``f'' con un tipo semejante, así que se devuelve error.

Como ya hemos mencionado, para seguir contrastando otros errores de tipado, se puede tomar cualquiera de los inputs 1, 2 y 3.
Estos son fragmentos de código más largos con múltiples errores de tipado. Todos los errores, como ya dijimos, están marcadores
en la línea del error por un comentario.